





\begin{theorem}[Printed Theorem C2]\label{mcrt-auto-theorem-c2}\dcref{C2}\source{matsumura-commutative-ring-theory:page-0299}\uses{}
\begin{verbatim}
Let A be a ring and M, N A-modules.                   Then

               i(M@N)=     @ [(iM)@,(iN)].
                          s+,=r
\end{verbatim}
\end{theorem}
\begin{proof}
\begin{verbatim}
@= l(M @ N) can be written as a direct sum of all possible r-fold
tensor products of copies of M and N (with 2? summands). Write L,,, for the
submodule which is the direct sum of all tensor products involving s copies
of M and t copies of N (with s + t = Y). Thus



Now let Q be the submodule of Ql(M ON)
generated by all elements of the form ... @x@...@x@...;           we have
Q = @[Q n L,,,]. We see at once that Q A L,,, is the submodule generated
by all elements of the forms ...@y@...@yO...         (with either HEM or
YEN), and aOyOfiOz@y+a@z@fi@y@y                    (with REM and ZEN).
Thus one sees easily that

          bW)=                &f@N)            /Q= @ &,,IQnLs,,),
                          (   1            )       S+t=T

and
          L,,,/L,,,nQ   -(i       M)O(   L W.     n
\end{verbatim}
\end{proof}
